problem:  
Let \( (M^m,g) \) be a compact Riemannian manifold with bounded geometry, and \( (N^n,h) \) a compact manifold with nonpositive sectional curvature.  
For a sequence of \( \varepsilon\)-approximately harmonic maps \(u_\varepsilon : M \to N\) satisfying  
\[
\|\tau(u_\varepsilon)\|_{L^2(M)} \le \varepsilon, \quad E(u_\varepsilon)\le E_0,
\]  
consider the stress–energy tensors  
\[
S_\varepsilon = \frac12 |\nabla u_\varepsilon|^2 g - (\nabla u_\varepsilon)^{\!*}\nabla u_\varepsilon,
\]
and the distributional divergence identity
\[
\operatorname{div} S_\varepsilon = -\langle \tau(u_\varepsilon),\nabla u_\varepsilon\rangle.
\]

Suppose \( S_\varepsilon \rightharpoonup S_* \) weakly as tensor-valued Radon measures. Let  
\(T := S_* - \tfrac12|\nabla u_*|^2 g + (\nabla u_*)^{\!*}\nabla u_*\)  
be the *tensorial defect measure*, representing the weak limit discrepancy in stress–energy beyond the scalar energy defect.

**Problem (revised focus):**  
At bubbling points, the scalar energy defect measure \(\nu\) captures lost energy through harmonic spheres.  Is it true that, after rescaling bubbles around such points,  
the anisotropic (trace-free) part of \(T\) satisfies a **quantitative isotropy principle**:

> There exists a dimension-dependent modulus \( \Gamma_m:(0,\infty)\to(0,\infty)\), with \(\Gamma_m(t)\to0\) as \(t\to0\), such that for every sequence of \( \varepsilon\)-approximately harmonic maps as above,
> \[
> \frac{\|T - \tfrac{1}{m}\operatorname{tr}_g T\cdot g\|(M)}{\nu(M)} \le \Gamma_m( \|\tau(u_\varepsilon)\|_{L^2(M)} ).
> \]
> 
> In other words, as the tension field becomes small, **the tensorial defect in stress–energy becomes asymptotically isotropic**, uniformly in the geometry class—a statement that would quantitatively rule out formation of anisotropic stress concentration at bubbling.

Equivalently: Does near-harmonicity enforce that, in all possible blow-up limits of \(u_\varepsilon\), the defect stress–energy tensor approaches a scalar multiple of the metric, to within a dimension-only modulus controlled by \( \|\tau(u_\varepsilon)\|_2 \)?  

Why is it a "good" problem:  
This refinement directly targets what standard bubbling theory cannot see: the *directional structure* of defect measures. Classical results describe only total scalar energy loss, not whether that lost energy is isotropically distributed in tangent directions.  

Asking whether small tension enforces asymptotic isotropy of stress–energy defects bridges several deep areas:

- From *analytic PDE*: one must analyze how the nearly divergence-free condition constrains the microlocal structure of tensor-valued Radon measures;
- From *geometric measure theory*: it introduces a new notion of “directional concentration rigidity,” extending scalar defect theory;
- From *geometric analysis*: it asks whether the harmonic map equation’s variational symmetry (stationarity under diffeomorphisms) quantitatively manifests in the isotropy of defect tensors.

A positive resolution would imply a new form of *quantitative stationarity rigidity*: not only does small tension bound total defect mass, it also enforces directional symmetry. Proving this would likely require new compensated compactness or H-measure tools linking weak convergence of quadratic forms to geometric isotropy.  

The statement is simple—“does small near-harmonicity force isotropy of stress–energy defects?”—but conceptually deep, highlighting subtle structure beyond classical energy quantization, and pointing towards a new layer of quantitative rigidity theory.